% Generated draft: requires mathematical and source review.
\documentclass[11pt]{article}
\usepackage[margin=25mm]{geometry}
\usepackage{fontspec}
\setmainfont{Latin Modern Roman}
\usepackage{amsmath,amssymb,mathtools,unicode-math}
\setmathfont{Latin Modern Math}
\usepackage{enumitem,needspace,xurl,hyperref}
\setlength{\parindent}{0pt}
\setlength{\parskip}{5pt}
\setlength{\emergencystretch}{3em}
\newcommand{\N}{\mathbb{N}}
\newcommand{\Z}{\mathbb{Z}}
\newcommand{\Q}{\mathbb{Q}}
\newcommand{\R}{\mathbb{R}}
\newcommand{\C}{\mathbb{C}}
\newcommand{\F}{\mathbb{F}}
\newcommand{\A}{\mathbb{A}}
\newcommand{\PP}{\mathbb{P}}
\newcommand{\E}{\mathbb{E}}
\newcommand{\eps}{\varepsilon}
\newcommand{\dd}{\,\mathrm{d}}
\begin{document}
\section*{UnsolvedMath ID: 30005415}
\textbf{Source problem code:} catalogue entry
\subsection{Definitions and mathematical statement}
Let \((R,\mathfrak m,k)\) be a complete Noetherian local Cohen--Macaulay normal domain, where \(k=R/\mathfrak m\), and put \(d=\dim R\). The characteristic is arbitrary. Let \(I\subsetneq R\) be an ideal of pure height \(h\). The meaning of pure height and the symbolic-power convention require the qualifications below.

Write \(\operatorname{Min}_R(R/I)\) for the primes minimal over \(I\), and \(\operatorname{Ass}_R(R/I)\) for the associated primes of \(R/I\). Two possible symbolic-power conventions are
\[
 I^{(n)}_{\min}=\bigcap_{\mathfrak p\in\operatorname{Min}_R(R/I)}(I^nR_{\mathfrak p}\cap R),\qquad
 I^{(n)}_{\mathrm{ass}}=\bigcap_{\mathfrak p\in\operatorname{Ass}_R(R/I)}(I^nR_{\mathfrak p}\cap R)
 \quad(n\geq1).
\]
The supplied record does not explicitly specify a convention. Likewise, pure height could require that every minimal prime over \(I\) have height \(h\), or could require that every associated prime of \(R/I\) have height \(h\). Under the latter requirement the two displayed symbolic-power conventions agree. These choices must be resolved before treating the draft as an unambiguous formulation.

For an \(R\)-module \(M\), let \(H^j_{\mathfrak m}(M)\) denote its \(j\)-th local cohomology module, the \(j\)-th right derived functor of \(\Gamma_{\mathfrak m}(M)=\{x\in M:\mathfrak m^t x=0\text{ for some }t\geq0\}\). An ideal \(\mathfrak a\) annihilates \(M\) if \(\mathfrak aM=0\). Its height is \(\operatorname{ht}(\mathfrak a)=\inf_{\mathfrak p\supseteq\mathfrak a}\operatorname{ht}(\mathfrak p)\).

\textbf{Question as stated in the supplied clean statement.} After fixing the intended pure-height and symbolic-power conventions, is it true that for every such \(R\) and \(I\), and every integer \(j\) with \(1\leq j\leq d-h-1\), there exists an ideal \(\mathfrak a_j\subseteq R\) such that
\[
 \operatorname{ht}(\mathfrak a_j)=d-j+1
 \quad\text{and}\quad
 \mathfrak a_j H^j_{\mathfrak m}(R/I^{(n)})=0
 \qquad\text{for every integer }n\geq1?
\]
Here \(\mathfrak a_j\) may depend on \(R,I,j\), but must be independent of \(n\). If \(d-h\leq1\), there are no indices \(j\) in the stated range. If the source includes zero in \(\N\), the convention \(I^{(0)}=R\) makes the additional case vacuous.

\textbf{Unresolved annihilation qualification.} The supplied original extraction contains the string \texttt{anj} in its annihilation equation. This may represent \(\mathfrak a_j^{\,n}\), in which case the intended condition would instead be
\[
 (\mathfrak a_j)^n H^j_{\mathfrak m}(R/I^{(n)})=0
 \qquad(n\geq1),
\]
with the same height and independence requirements. The clean statement asks for annihilation by \(\mathfrak a_j\) itself. The supplied verification label does not resolve this discrepancy within the evidence provided here.

\textbf{Supplied contextual claim.} The original passage describes the problem as characteristic-free and says that solving it would settle the weak-implies-strong conjecture. It further says that solving it for ideals of pure height \(1\) in a weakly \(F\)-regular ring would suffice to settle its Conjecture~1. These are retained as source claims; the record does not supply Conjecture~1 or its precise hypotheses. Weak \(F\)-regularity is not an additional hypothesis of the general question above.
\subsection{Short English statement}
The supplied clean statement asks whether each indicated local cohomology degree admits an ideal of exactly the prescribed height that annihilates the modules for all symbolic powers of \(I\), using one ideal independent of the symbolic-power index. The original extraction may instead ask for annihilation by the corresponding power of that ideal. The annihilation equation, pure-height convention, and symbolic-power convention require human review. The record labels the problem open; that status is supplied dated curation, not an independently verified conclusion.
\subsection{Sources}
\begin{itemize}
\item Local registry: supplied UnsolvedMath record, permanent ID \(30005415\), title \emph{Uniform Local Cohomology Annihilators for Symbolic Powers}, exact source problem code \texttt{catalogue entry}, in the Oberwolfach Reports Open Problems collection. No individual proposer is identified in the record; the supplied proposed year is 2023.
\item Supplied report citation: \emph{Resolutions in Local Algebra and Singularity Theory} (2023), Oberwolfach Reports, pp.~327--395, DOI \texttt{10.4171/owr/2023/6}. Source URL and literature-source label ``original OWR open problem'': \url{https://doi.org/10.4171/owr/2023/6}. No specific problem page within the report is supplied.
\item Supplied literature source labeled ``partial weak-implies-strong results'': \url{https://par.nsf.gov/servlets/purl/10501870}. The record supplies no author, title, or precise theorem from this reference.
\item Supplied literature assessment, checked on \textbf{2026-08-22}: the characteristic-free annihilation problem is described as open, with no general solution known, and its solution is said to imply weak-implies-strong. The background describes this as dated literature triage derived from the supplied OWR curation and requests independent verification of time-sensitive status and cited claims. This is supplied dated curation, not independently verified here.
\item Supplied statement-verification metadata: the label is \texttt{corrected\_verified}, and the record reports that the official Oberwolfach report was checked directly and that the clean statement restores definitions and removes extraction artifacts. This reports the registry's verification claim; it does not establish independent verification in this draft or resolve the ambiguous original annihilation expression.
\end{itemize}
\end{document}
